% Requires amsmath, amssymb, and the paper's existing citation package.
% Uses the citation keys in Sam's bibliography.
% Replace the older subsection carrying app:comparison-bounds; do not include both.
% The final paragraph states the uniform estimates required from the QUILLS proof.

\subsection{Derivation of the comparison bounds}
\label{app:comparison-bounds}

\paragraph{Common assumptions.}
Fix a uniformly bounded class of functions analytic on a fixed complex
neighborhood of $[-1,1]$. Constants below may depend on that class, but not
on the requested error $\varepsilon$. Error is measured in
$\|g\|_{2,\Omega}=(\tfrac12\int_{-1}^{1}|g(x)|^2\,dx)^{1/2}$;
the uniform error bounds for the competing constructions imply this bound.
Let $W$ denote maximum hidden-layer width, $B$ the largest absolute weight
or bias, and $b$ the number of working significand bits. Assume target
samples with absolute error $O(u)$, accurately rounded elementary functions,
and sufficient exponent range, where $u\asymp2^{-b}$ is unit roundoff.
If construction and evaluation add at most $uA$ error, then
\begin{equation}
 b\ge\log_2(1/\varepsilon)+\log_2\max\{1,A\}+O(1)
 \label{eq:table-precision}
\end{equation}
suffices. We allocate fixed fractions of the error budget to approximation
and arithmetic. All bounds concern specified implementations; none is a
lower bound on necessary precision. Here $\widetilde O$ suppresses powers
of $\log\log(1/\varepsilon)$.

\paragraph{Classical staircase.}
For $x_j=-1+2j/N$ and $S(t)=(1+\tanh t)/2$, take
\[
 q_N(x)=f(-1)+\sum_{j=1}^{N}
 [f(x_j)-f(x_{j-1})]S(N(x-x_j)).
\]
The analytic class has a uniform Lipschitz bound, so each coefficient
increment is $O(N^{-1})$. The piecewise-constant approximation error is
$O(N^{-1})$; smoothing its jumps adds the same order, since the sigmoid
tails decay exponentially with distance in grid cells. Hence
$\|f-q_N\|_\infty=O(N^{-1})$ and $N=O(\varepsilon^{-1})$ suffices.
There are $N$ tanh neurons; slopes and biases are $O(N)$.
Coefficient formation and direct evaluation have error $O(uN)$, using
the bounded total absolute readout weight. Thus
\[
 W=O(\varepsilon^{-1}),\qquad B=O(\varepsilon^{-1}),\qquad
 b=O(\log(1/\varepsilon)).
\]
Jones~\citeyearpar{58342} is the historical reference; these estimates
are for the explicit staircase above.

\paragraph{Mhaskar: an explicit tanh specialization.}
The polynomial construction underlying~\cite{mhaskar1996approximation}
can be realized as follows. Analyticity gives a degree-$n$ polynomial
with error $C\rho^{-n}$, for fixed $\rho>1$, so
$n=O(\log(1/\varepsilon))$ suffices. For $1\le p\le n$, replace $x^p$ by
\[
 H_{p,h}(x)=
 \frac{\sum_{r=0}^{p}(-1)^{p-r}\binom pr
 \tanh\!\left(\beta+(r-p/2)hx\right)}
 {h^p\tanh^{(p)}(\beta)},\qquad \beta=\tfrac12\log2,
\]
where $h>0$ is the finite-difference spacing. All these formulas share
the $2n+1$ slopes $hj/2$, $-n\le j\le n$, so one hidden layer of
width $O(n)$ suffices.

The denominators are quantitatively nonzero. Indeed,
$\tanh^{(p)}(\beta)=2^{p+2}E_p(-2)/3^{p+1}$, where
$E_1(z)=1$ and
$E_{p+1}(z)=(1+pz)E_p(z)+z(1-z)E_p'(z)$.
These integer polynomials have constant term one; hence $E_p(-2)$
is odd and $|\tanh^{(p)}(\beta)|^{-1}\le e^{Cp}$.
The centered-difference error satisfies
\[
 \|H_{p,h}-x^p\|_\infty
 \le \frac{ph^2\|\tanh^{(p+2)}\|_\infty}
 {24|\tanh^{(p)}(\beta)|}.
\]
Together with Cauchy bounds on tanh derivatives and the
$e^{O(n)}$ cost of converting a bounded Chebyshev expansion to monomials,
this bounds the combined error by $h^2e^{Cn\log(n+2)}$.
Choose dyadic $h$ small enough that this is $O(\varepsilon)$, with
$\log(1/h)=O(\log(1/\varepsilon)+n\log(n+2))$.
Summing absolute finite-difference coefficients gives
$B\le e^{Cn}(2/h)^n$.
Compute the $E_p(-2)$ numerators with exact integer arithmetic and use
stable Chebyshev coefficient computation; coefficient formation and
network evaluation then have amplification at most
$A\le e^{Cn\log(n+2)}(2/h)^n$.
Consequently,
\[
 W=O(\log(1/\varepsilon)),\qquad
 \log B,\ \log A
 =O\!\left(\log^2(1/\varepsilon)\log\log(1/\varepsilon)\right).
\]
Equation~\eqref{eq:table-precision} gives the stated
$\widetilde O(\log^2(1/\varepsilon))$ precision bound.
These parameter and arithmetic bounds are derived for this specialization.

\paragraph{Costarelli--Spigler: sampled tanh implementation.}
Section~5 of~\cite{Costarelli2015Approximation} uses
$\sigma(t)=(1+\tanh t)/2$, $\phi=\sigma'$, and
$\psi(t)=\sigma(t+1)-\sigma(t)$. The translates of $\psi$ are nonnegative
and sum to one. Extend $f$ to a bounded Lipschitz function $F$ by linearly
tapering its endpoint values to zero on $[-2,-1]$ and $[1,2]$, and setting
$F=0$ elsewhere outside $[-2,2]$.
Their Theorem~5.4, with scale $N$ and truncation $|k|\le2N$, gives
$O(N^{-1})$ error using coefficients
\[
 a_k=\int_{\mathbb R}\phi(t)F((k+t)/N)\,dt.
\]
If $L_F$ is a Lipschitz constant for $F$, then
\[
 |a_k-F(k/N)|\le\frac{L_F}{N}
 \int_{\mathbb R}|t|\phi(t)\,dt=O(N^{-1}).
\]
The partition of unity implies that replacing every $a_k$ by the sample
$F(k/N)$ adds only $O(N^{-1})$ error. Thus the explicit sampled network
\[
 q_N(x)=\sum_{k=-2N}^{2N}F(k/N)\psi(Nx-k)
\]
retains that rate without quadrature. Expanding each $\psi$ into two tanh
terms gives one hidden layer with $W=O(N)$, slopes and biases $O(N)$,
and bounded readout coefficients. A conservative bound for sample-based
construction and direct evaluation is $O(uN^2)$. Taking
$N=O(\varepsilon^{-1})$ and applying~\eqref{eq:table-precision} gives
\[
 W=O(\varepsilon^{-1}),\qquad B=O(\varepsilon^{-1}),\qquad
 b=O(\log(1/\varepsilon)).
\]
The sampled implementation and its precision bound are additional
derivations, not bit-complexity claims made in that paper.

\paragraph{ChebNet.}
For a degree-$n$ Chebyshev polynomial,
Theorem~1 of~\cite{tang2019chebnet} gives an exact ReQU realization
with $O(n)$ neurons and $O(\log n)$ hidden layers.
Analyticity supplies coefficients $|c_j|\le C\rho^{-j}$ and requires
$n=O(\log(1/\varepsilon))$.
In the recursive construction, an original coefficient $c_j$ is
multiplied by at most $2\max\{1,j\}$ along any branch: every doubling
reduces its index by at least half. Therefore the coefficient sums
of intermediate polynomials are bounded by
$2\sum_j\max\{1,j\}|c_j|=O(1)$.
The ReQU implementation has bounded weights and intermediates.
Bounded fan-in and $O(\log n)$ depth give at most polynomial-in-$n$
arithmetic amplification; stable coefficient construction does also.
Thus
\[
 W=O(\log(1/\varepsilon)),\qquad B=O(1),\qquad
 b=\log_2(1/\varepsilon)+O(\log\log(1/\varepsilon)).
\]
The hidden-layer count is $O(\log\log(1/\varepsilon))$.
The weight and precision estimates are derived for this analytic class.

\paragraph{QUILLS: estimates required from the main theorem.}
The QUILLS row requires the main construction and recovery result to
provide, along the bandwidth choice used to attain error $\varepsilon$,
\[
 W=O(\log(1/\varepsilon)),\qquad
 \|c\|_1=O(1),\qquad
 \text{construction and evaluation error}\le u\,\mathrm{poly}(W),
\]
with constants independent of $\varepsilon$; $c$ is the physical readout
vector. With spacing $h\asymp W^{-1}$, bounded center locations $x_j$,
and bounded relative bandwidth $\lambda$, the hidden slopes and biases
are $\gamma=\lambda/h=O(W)$ and $-\gamma x_j=O(W)$.
These estimates give one tanh hidden layer and
\[
 B=O(\log(1/\varepsilon)),\qquad
 b=\log_2(1/\varepsilon)+O(\log W)=O(\log(1/\varepsilon)).
\]
The uniform bandwidth and recovery estimates are inputs from the main
proof, not consequences of logarithmic width alone. In particular,
constants established only for fixed $\lambda$ must be controlled before
allowing $\lambda$ to vary with $\varepsilon$.